\documentclass[11pt]{article}

\usepackage[margin=1in]{geometry}
\usepackage[utf8]{inputenc}
\usepackage[T1]{fontenc}
\usepackage{amsmath,amssymb}
\usepackage{booktabs}
\usepackage{array}
\usepackage{graphicx}
\usepackage{xcolor}
\usepackage{enumitem}
\usepackage{caption}
\usepackage{longtable}
\usepackage{float}
\usepackage{algorithm}
\usepackage{algpseudocode}
\usepackage[hidelinks,colorlinks=false]{hyperref}
\usepackage[noabbrev,capitalise]{cleveref}

\setlength{\parindent}{0pt}
\setlength{\parskip}{6pt}
\emergencystretch=3em

\definecolor{cIndigo}{HTML}{3A4CC4}
\definecolor{cAmber}{HTML}{96590A}
\definecolor{cInk2}{HTML}{525A6C}
\definecolor{cPanel}{HTML}{F6F7F9}
\definecolor{cRule}{HTML}{BFC6D3}

\newcommand{\dgraph}{d_{\mathrm{graph}}}
\newcommand{\zv}{\mathbf{z}}

\title{\vspace{-2em}\textbf{Predictor-Based Cross-Episode Stitching}\\[0.3em]
\large Using LeWM's action-conditioned predictor to verify and discover latent-graph edges,
and the case for testing synthetic HER before trusting it}
\author{Companion proposal to \texttt{main.tex}}
\date{6 September 2026}

\begin{document}
\maketitle
\vspace{-1.5em}

\section{What this proposes, in one paragraph}

The latent graph in \texttt{main.tex} identifies cross-episode edges by a single test: are two
landmarks' encoded states within a calibrated Euclidean threshold of each other. That test only
ever looks at where a trajectory \emph{already is}, so it can never propose an edge between two
states that are dynamically one action apart but happen to sit at different points in latent
space --- a slightly different path to the same doorway, for instance. LeWM's own
action-conditioned predictor $\mathrm{predict}(\zv,\mathbf{a})$ is differentiable and frozen, and
is exactly the tool needed to ask a stronger question than metric proximity: is there an action
that actually connects these two states. The current methodology answers that with one
mechanism: a small inverse-dynamics-model (IDM) ensemble that infers the connecting action for
a candidate pair directly, with the frozen predictor as an independent cross-check
(\cref{sec:method}) --- and, critically, a staged experiment ladder (\cref{sec:ladder}) for
testing it \emph{before} trusting the riskiest downstream idea: replacing the graph's existing,
already-validated auxiliary-regression mechanism with synthetic cross-episode HER tuples built
from these predicted actions. \Cref{sec:risk} argues that idea has a real chance of hurting
performance, not just failing to help, and \cref{sec:e1-results} reports the first experiment's
result.

\section{Method}
\label{sec:method}

\textbf{IDM ensemble confidence, verified against the predictor.} This is the only mechanism
actually in use. Train $N{=}5\text{--}8$ small MLP heads
$\mathrm{IDM}_\theta(\zv_1,\zv_2) \to \hat a$, supervised on real adjacent
$(\zv_t,\zv_{t+1},a_t)$ triples --- a genuinely different model from the predictor above,
since it inverts the dynamics rather than rolling them forward. Given a candidate cross-episode
pair $(\zv_i,\zv_j)$, the ensemble infers the action each member believes would connect them;
$\mathrm{Var}(\hat a)$ across the ensemble is a confidence signal, since independently-trained
regressors should disagree more on pairs unlike anything in training. The caveat: some of that
variance may reflect genuine action multimodality (several different actions produce a similar
effect) rather than lack of support, so variance alone is not trusted as a confidence measure ---
it is combined with a second, independent check: the ensemble's mean action is fed through the
frozen predictor, and the pair is accepted only if the resulting prediction actually reproduces
$\zv_j$. Requiring agreement between two separately-trained models (the IDM, supervised
end-to-end on real transitions; the predictor, pretrained on a different objective entirely) is
a strictly stronger claim than either model's self-consistency alone. \Cref{sec:e2-results}
gives the exact construction and both verification stages (a latent-only check, then a
ground-truth environment check) in full.

\textit{Supporting latent-space diagnostics --- consecutive vs.\ random-pair distances, a
theoretical Gaussian reference, predictor error and direction-agreement, the embedding's PCA
spectrum, IDM ensemble variance across three populations, cosine similarity, and (as of a
later correction) a cross-episode splice verification --- are collected separately in
\texttt{docs/graph-proposal/latent-diagnostics.tex}, generated by
\texttt{scripts/latent\_diagnostics.py}. \textbf{Important correction:} every predictor call
below this point in the present document, and every predictor-error/direction number
originally reported in that companion document, used a WRONG action invocation --- a single
real action zero-padded into the 10-dim slot $\mathrm{action\_encoder}$ expects, compared
one real step later. The checkpoint's actual convention (\texttt{docs/lewm.txt} App.\ D) is a
frame-skip of $5$: one predictor call consumes $5$ concatenated real actions
($5\times2{=}10$ dims, no padding) and predicts $5$ real steps ahead. The companion document
now reports corrected numbers throughout; $\tau_{\mathrm{fwd}}$'s mis-scaling below is
diagnosed using the (still-valid) general argument, but its specific value and the E2
stage-1/3 experiments that followed were computed under the wrong invocation and have not
been re-run correctly. What IS now validated under the corrected invocation: the IDM
variance finding is unaffected (it never touched the predictor) --- real transitions and the
candidate pool overlap almost completely, genuinely random pairs sit an order of magnitude
higher ($0.0495$ vs.\ $\sim\!0.005$); and, most importantly, a corrected cross-episode splice
test (feed the IDM's inferred action into a real 5-action block, continue with the target
episode's own next 4 actions, predict from $\zv_i$) shows the IDM's action performing
\emph{statistically indistinguishably} from a real action on Two-Room (hybrid mean error
$5.59$ vs.\ real-block reference $5.77$) --- direct, positive evidence for the underlying
mechanism that the untested "candidate for a fix" framing below could not yet claim.}

\section{The central risk: synthetic HER is not free}
\label{sec:risk}

Standard same-episode hindsight relabeling is, in a precise sense, free training data: the
relabeled goal is defined as whatever state the trajectory actually reached, so the label is
trivially true by construction and carries no noise. The mechanism in \cref{sec:method} that
identifies a \emph{cross-episode} connection breaks that property the moment its output ---
an inferred action --- is treated as a real transition. If the inference is wrong, the
resulting $(s,a,s',g)$ tuple is a genuine falsehood about the environment, not a relabeling of
something that happened.

That distinction matters most for where the tuple is used. Feeding a graph-derived distance into
the existing auxiliary regression term on $V(s,g)$ (\texttt{main.tex}'s mode 2 consumption
algorithm) injects a soft, scalar bias that is still tempered by $V$'s own TD/expectile regression
against real Bellman targets --- and even that soft channel has already been shown, on Push-T's
noisy tiers, to invert the actor's data preference under the wrong scale (see
\texttt{main.tex}, Push-T section). Feeding a synthetic action directly into the actor's
imitation loss, $(\hat a - a_{\mathrm{pred}})^2$, has no such anchor: a wrong target is a wrong
target, with nothing pulling the policy back toward grounded behavior.

There is a direct precedent in this project for why a mechanism that looks safe on paper can
still fail in practice: potential-based reward shaping (\texttt{main.tex}'s mode 3) is
\emph{provably} policy-invariant --- a strictly stronger theoretical safety guarantee than
anything available for synthetic HER --- and it still lost on 4 of 5 Two-Room tiers, because TD
compounds the graph's noise even though the shaping itself cannot change the optimal policy in
the exact case. Synthetic cross-episode HER has a weaker guarantee than mode 3 had. There is
therefore no basis to assume the gap between a baseline and a synthetic-HER-augmented policy stays
non-negative, and every reason to expect the false-edge count --- and thus the risk --- to grow,
not shrink, as the graph gets larger and the number of candidate cross-episode pairs grows with
it. \Cref{sec:ladder} exists to measure this rather than assume an answer either way.

\section{Experiment ladder}
\label{sec:ladder}

Six gated stages, Two-Room first (a trustworthy oracle distance exists to diagnose against),
Push-T only once the mechanism is trusted --- the same sequencing \texttt{main.tex} used for its
own B0--B5 ladder. Each stage's kill criterion is checked before the next stage is attempted.

\begin{center}
\small
\begin{longtable}{@{}p{0.6cm}p{4.6cm}p{5.4cm}p{3.6cm}@{}}
\caption{Experiment ladder for predictor-based cross-episode stitching.} \\
\toprule
& \textbf{Question} & \textbf{Method} & \textbf{Kill criterion / status} \\
\midrule
\endfirsthead
\multicolumn{4}{l}{\small\itshape (Experiment ladder, continued)} \\
\toprule
& \textbf{Question} & \textbf{Method} & \textbf{Kill criterion / status} \\
\midrule
\endhead
\bottomrule
\multicolumn{4}{r}{\small\itshape (continued on next page)} \\
\endfoot
\bottomrule
\endlastfoot
E1 & Can an IDM ensemble recover a known action at all? & Train on real adjacent Two-Room
transitions; test on held-out same-episode pairs (\cref{alg:idm} sketch). & \textcolor{cAmber}
{\textbf{Done --- passed.}} 98.4\% MSE reduction vs.\ marginal baseline
(\cref{sec:e1-results}). \\
\addlinespace
E2 & Does IDM-ensemble variance track true confidence, and does a wide kNN candidate pool,
filtered by ensemble variance and verified against the frozen predictor, beat plain kNN on
precision? & Wide top-20 kNN candidates on Two-Room's B0 landmark set; IDM ensemble infers
the connecting action per pair; accept iff variance is typical AND the predictor's forward
step reproduces the target. & \textcolor{cAmber}{\textbf{Open.}} The absolute-floor criterion
FAILED to discriminate; ground-truth env stepping shows the IDM's inferred actions are
genuinely correct, even cross-trajectory, so the mechanism itself is sound
(\cref{sec:e2-results}). A figure-motivated fix (recalibrated variance $+$ direction
agreement) was tried next and made precision \emph{worse} ($0.717\to0.609$), for a diagnosed,
scale-dependent reason --- no working latent-only criterion exists yet. E3 remains blocked
on this. \\
\addlinespace
E3 & Does graph quality from these edges improve, or degrade, with scale? & Re-run the B0
Spearman-vs-oracle gate at 50/300/1000-episode landmark counts with E2's edges added. & Not yet
run. Kill criterion: no improvement, or degradation at larger scale, answers the scaling
question directly. \\
\addlinespace
E4 & Design A --- do verified edges help when fed only to the existing (safe) aux-regression
mechanism, no synthetic actor targets? & Wire E3's edges into $\phi_{\mathrm{dist}}$ as ordinary
identification edges; sweep tiers $\times$ 3 seeds. & Not yet run. Kill criterion: if this
doesn't help, Design B is even less likely to. \\
\addlinespace
E5 & Design B --- does synthetic cross-episode HER, fed into $\mathrm{actor\_loss}$, ever beat
baseline, and does the gap survive scale? & Sweep a confidence-gated synthetic-tuple fraction
$\in\{0, 0.05, 0.15, 0.30, 1.0\}$ across tier sizes. & Not yet run. This is the concrete test of
\cref{sec:risk}'s concern: a negative gap at any fraction/scale confirms it. \\
\addlinespace
E6 & Does the winning Two-Room configuration transfer to Push-T? & Repeat E4/E5's best config
under \texttt{eval\_protocol=same\_episode}; watch for the known expert-vs-noisy AWR
weight-inversion failure mode. & Not attempted before Two-Room shows a clear non-negative
signal. \\
\end{longtable}
\end{center}

\section{Results}

\subsection{E1: IDM ensemble sanity check}
\label{sec:e1-results}

\begin{algorithm}[H]
\caption{E1 methodology}
\label{alg:idm}
\begin{algorithmic}[1]
\State \textbf{Data:} Two-Room's existing B0 landmark cache, 300 episodes $\to$ 27{,}675 frames
\State \textbf{Pairs:} all real adjacent (transition) pairs, 27{,}375 total
\State \textbf{Split:} 80/20 by \emph{episode} (240 train / 60 test episodes $\to$ 21{,}910 /
5{,}465 pairs) --- same convention as the B3 tier-setup split
\State \textbf{Model:} 6 independent MLPs (hidden 256), each on its own bootstrap resample and
weight init, regressing $a_i$ from $\mathrm{concat}(\zv_i,\zv_j)$, 3000 Adam steps
\State \textbf{Baseline:} predict the training-set mean action for every test pair (the
``marginal action distribution'')
\State \textbf{Kill criterion:} held-out ensemble MSE must be meaningfully below the marginal
baseline's MSE
\end{algorithmic}
\end{algorithm}

\textbf{Results.}

\begin{table}[h]
\centering
\caption{E1 result. MSE on held-out same-episode transitions, 60 held-out episodes never seen
in training.}
\begin{tabular}{lc}
\toprule
& MSE \\
\midrule
marginal-action baseline & 0.7544 \\
per-member average & 0.0170 \\
ensemble mean & \textbf{0.0120} \\
\midrule
improvement over baseline & \textbf{98.4\%} \\
\bottomrule
\end{tabular}
\end{table}

Ensemble variance on these (in-distribution) held-out pairs: mean $0.0051$, std $0.0075$.

\textbf{Takeaways.}
\begin{itemize}[leftmargin=1.4em,itemsep=2pt]
\item The ensemble reliably inverts Two-Room's one-step dynamics on pairs shaped like its
training distribution --- real, same-episode, one step apart. The kill criterion (a clear
improvement over predicting the mean action) passes by a wide margin, not a marginal one.
\item This says nothing yet about genuinely cross-episode pairs, which is what E2 tests.
Variance being small and reasonably stable here is the necessary baseline for E2 to compare
against, not evidence that variance will behave the same way on out-of-distribution pairs.
\item This clears the door to E2. It does not, on its own, provide any evidence about
\cref{sec:risk}'s actual concern --- that is a claim about E5, five stages away.
\end{itemize}

\subsection{E2: kNN candidates, IDM variance filter, predictor verification}
\label{sec:e2-results}

Rather than borrowing an unrelated real action from a neighboring state (the earlier,
abandoned design described in \cref{sec:method}), have the IDM ensemble infer the action
connecting a SPECIFIC candidate pair $(\zv_i,\zv_j)$ directly, and accept the pair only if
two independently-trained models agree: the ensemble's own variance is no worse than
typical, \emph{and} feeding its mean action through the frozen predictor reproduces $\zv_j$.
Tested in two stages --- a cheap latent-only check first, then a ground-truth check in the
real environment on whatever the latent check couldn't settle.

\begin{algorithm}[H]
\caption{E2 stage 1: candidate construction and latent-space verification}
\label{alg:e2-latent}
\begin{algorithmic}[1]
\State \textbf{Candidates:} unfiltered top-20 nearest cross-episode neighbors by raw latent
distance (\texttt{raw\_topk\_cross\_episode\_pairs}), NOT the calibrated $k{=}4$
identification-edge threshold --- 298{,}983 pairs, chosen deliberately looser so the
verification step has a real mix of true/false candidates to discriminate, unlike testing it
on an already-92.7\%-precision set
\State \textbf{IDM ensemble:} 6 MLPs, same architecture and training split as E1, trained on
real adjacent transitions only
\For{each candidate pair $(\zv_i,\zv_j)$}
    \State $\{\hat a_1,\dots,\hat a_6\} \gets \{\mathrm{IDM}_k(\zv_i,\zv_j)\}$; $\bar a \gets$
    their mean; $\mathrm{Var}(\hat a) \gets$ their variance
    \State \textbf{require} $\mathrm{Var}(\hat a) < \tau_{\mathrm{var}}$ \Comment{$\tau_{\mathrm{var}} = $
    90th percentile of the ensemble's OWN variance on real held-out adjacent pairs --- ``no
    worse than typical''}
    \State $\hat\zv \gets \mathrm{predict}(\zv_i, \mathrm{enc}_a(\bar a))$
    \State \textbf{require} $\|\hat\zv - \zv_j\|^2 < \tau_{\mathrm{fwd}}$ \Comment{$\tau_{\mathrm{fwd}} = $
    90th-percentile one-step predictor error on real transitions, 51.0}
    \State \textbf{accept} iff both requirements hold
\EndFor
\end{algorithmic}
\end{algorithm}

\textbf{Stage 1 result: does not discriminate.} 87.8\% of all candidates pass both checks;
precision (oracle distance $\le 2$) moves from 0.711 unfiltered to only 0.719 filtered;
recall on pairs already known to be good (0.876) is statistically indistinguishable from the
overall pass rate (0.878) --- the filter accepts or rejects almost independently of
correctness.

\textbf{Diagnosis.} Candidate pairs sit only $2.74$ apart in raw latent distance on average;
a real one-step transition moves $11.76$ on average --- a $4.3\times$ gap. $\tau_{\mathrm{fwd}}$
was calibrated at the larger (real-transition) scale, so a candidate pair already well inside
it before any action is applied passes almost automatically. Confirmed directly: the IDM's
inferred mean-action norm on these candidates is $0.42$ versus $1.20$ for real training
actions --- the ensemble is agreeing on a near-zero, near-identity ``explanation'' that
trivially satisfies both checks without asserting anything dynamically real. This could be a
fixable calibration-scale mismatch (e.g.\ a RELATIVE-improvement criterion instead of an
absolute floor borrowed from a different scale) rather than evidence the two-model idea is
unsound --- Stage 2 settles which.

\begin{algorithm}[H]
\caption{E2 stage 2: ground-truth environment verification}
\label{alg:e2-env}
\begin{algorithmic}[1]
\Require candidate $(i,j)$; its inferred mean action $\bar a$ from stage 1
\State reset the real Two-Room environment at $\mathrm{state}_i$, with
$\mathrm{options}=\{\text{state}: \mathrm{state}_i,\ \text{target\_state}: \mathrm{state}_j\}$
\State force $\mathrm{speed} \gets 5.0$ \Comment{the dataset's empirically fixed value ---
mean 4.989, std 0.147, capped at exactly 5.0 over 200 episodes' real unclamped transitions
--- overriding \texttt{reset()}'s otherwise-randomized per-episode draw, which would
otherwise confound the replay}
\State step the environment with $\bar a$; read the REAL resulting position
$\mathrm{state}_i'$
\State compare $\|\mathrm{state}_i' - \mathrm{state}_j\|$ against: (a) $\|\mathrm{state}_i -
\mathrm{state}_j\|$ (a no-op control), and (b) the same step under a random action of equal
magnitude (a control for ``any small nudge helps'')
\end{algorithmic}
\end{algorithm}

\textbf{Stage 2, result 1 (all top-20 candidates): the action is real, not degenerate.} Mean
real distance $1.66 \to 0.35$ after the action ($79\%$ reduction), beating both the no-op
control ($1.66$, identical by construction) and the random-direction control ($2.30$, worse
than the no-op, as expected for pure noise) on $95.4\%$ / $95.6\%$ of pairs respectively. But
these top-20 pairs are overwhelmingly near-duplicates in real position (only $0.8\%$ of the
whole 298{,}983-pair pool sit more than $10$ apart), so this alone only shows local
correction between already-similar states.

\textbf{Stage 2, result 2 (restricted to the genuinely far-apart $4.2\%$ of the pool, real
distance $\ge 5$): real cross-trajectory stitching.} Mean inferred-action norm here is
$1.14$ --- essentially matching real training actions' $1.20$, not a degenerate near-zero
action, because the task genuinely requires real movement this time. Median real distance
$6.87 \to 1.69$ ($75\%$ reduction) after the action; $95.6\%$ of pairs improve over doing
nothing, $91.8\%$ improve over the random-direction control. The environment's own success
threshold (distance $<16.0$) barely moves ($0.784 \to 0.796$), because many of these pairs
are simply farther apart than one action (bounded by $\mathrm{speed}=5$) can close in a
single step regardless of correctness --- the median-distance reduction and the
beat-the-random-control margin are the informative numbers, not the raw hit rate against a
threshold tuned for full-episode navigation.

\textbf{Takeaways.}
\begin{itemize}[leftmargin=1.4em,itemsep=2pt]
\item Stage 2 revises stage 1's verdict. The IDM's inference mechanism --- ask it directly
what action connects a specific candidate pair --- is genuinely good, confirmed in the real
environment for pairs that are actually far apart, not merely near-duplicates. What failed in
stage 1 was its LATENT acceptance criterion (an absolute floor calibrated at the wrong length
scale), not the underlying inference.
\item Practical implication: do not run real-environment stepping as the production
verification mechanism at scale or for a new environment without a clean resettable oracle
(Push-T) --- that defeats the point of a learned latent-space graph, and Two-Room's fast grid
oracle is a lucky exception, not the general case. Use stage 2's ground truth as a
CALIBRATION set instead: fix stage 1's acceptance criterion (the relative-improvement fix
proposed above) against what real stepping confirms here, then trust the resulting cheap,
latent-only proxy for construction at scale.
\item This reopens the door to E3 onward, on a corrected latent criterion --- attempted next,
in stage 3.
\end{itemize}

\textbf{Stage 3: the corrected criterion, tried and failed differently.} Two fixes were
implemented, motivated directly by \texttt{latent-diagnostics.tex}'s figures: the variance
threshold recalibrated as the geometric mean of the real-transition and random-pair variance
medians (Figure 6: real $0.00505$, candidate pool $0.00560$, random $0.0495$), and the
absolute forward-error floor replaced with direction-of-displacement agreement,
$\cos(\mathrm{predict}(\zv_i,\bar a)-\zv_i,\ \zv_j-\zv_i) > 0.5$ (Figure 4's median $0.924$ on
real transitions). Re-run on the same candidate pool, this made things \emph{worse}, not
better: precision $0.717 \to 0.609$, recall on known-good pairs only $46.5\%$.

A distance-bucket breakdown pins down why:

\begin{table}[h]
\centering
\small
\caption{Candidates bucketed by raw latent distance $\|\zv_i-\zv_j\|$ (3000-pair oracle-scored
sample).}
\begin{tabular}{lcccc}
\toprule
raw distance & frac.\ actually good & mean $\cos_{\mathrm{dir}}$ & recall (good kept) &
precision among kept \\
\midrule
$[0,1)$  & $0.994$ & $0.144$ & $0.292$ & $1.000$ \\
$[1,2)$  & $0.947$ & $0.292$ & $0.377$ & $0.950$ \\
$[2,3)$  & $0.777$ & $0.395$ & $0.463$ & $0.780$ \\
$[3,5)$  & $0.489$ & $0.635$ & $0.691$ & $0.445$ \\
$[5,10)$ & $0.077$ & $0.805$ & $1.000$ & $0.085$ \\
\bottomrule
\end{tabular}
\end{table}

Two distinct failures, at opposite ends of the scale. Near-duplicate pairs (the bulk of the
pool, the first two rows) are $95$--$99\%$ genuinely good, but a target displacement this
tiny has a direction dominated by encoder noise, not real geometry --- mean cosine is only
$0.14$--$0.29$ for pairs that are almost all correct, so the filter wrongly rejects most of
the pool's best pairs. This is the same scale-mismatch failure as the original absolute
floor, relocated to a different statistic rather than fixed. At larger gaps (last row), only
$7.7\%$ are actually good yet mean cosine is $0.805$ and everything kept is $91.5\%$ wrong ---
the predictor's displacement direction does not track the specific target at this scale, it
looks like a generic directional bias rather than genuine confirmation.

\textbf{This is not a threshold-tuning problem.} No single global cosine threshold fixes both
failure modes at once: the small-gap failure needs a lower bar (or no direction check at
all), while the large-gap failure needs a direction signal that is target-specific, which a
single dot product can't provide regardless of where the bar is set. A working fix would need
to be scale-aware --- e.g.\ skip the model-based check entirely below some raw distance, since
near-duplicates are already well-served by the existing metric threshold alone --- and would
still need to solve the large-gap problem separately; neither is attempted here.

\section{Status}

E1 complete and passed. E2 (wide kNN candidates, IDM variance filter, predictor
verification) is \textbf{not complete} after three attempts (\cref{sec:e2-results}): stage
1's original absolute-floor criterion failed to discriminate; stage 2's ground-truth
environment stepping confirmed the underlying mechanism --- the IDM directly inferring the
connecting action for a candidate pair --- is genuinely correct, including for real
cross-trajectory pairs; but stage 3's figure-motivated fix (recalibrated variance threshold
plus direction-of-displacement agreement in place of the absolute floor) made precision
\emph{worse} ($0.717\to0.609$), for a diagnosed, scale-dependent reason (direction is
undefined for near-duplicates and non-specific at larger gaps) rather than a tuning mistake.
\textbf{Net effect:} the mechanism is validated, but no working latent-only acceptance
criterion exists yet --- E3 onward remains blocked until one is found, most likely a
scale-aware criterion that treats near-duplicate and larger-gap candidates differently rather
than a single global threshold of any kind. Full artifacts: \texttt{scripts/idm\_gate.py},
\texttt{scripts/common/idm\_ensemble.py}, \texttt{config/graph/idm\_gate.yaml},
\texttt{scripts/cycle\_consistency\_gate.py},
\texttt{config/graph/cycle\_consistency\_gate.yaml},
\texttt{scripts/env\_verify\_cycle\_edges.py}, results at
\texttt{outputs/idm\_gate\_tworoom/e1\_idm\_gate\_results.json},
\texttt{outputs/idm\_gate\_tworoom/e2b\_cycle\_consistency\_gate\_results.json}, and
\texttt{outputs/idm\_gate\_tworoom/e2c\_env\_verify\_results.json}. The staged plan this
document
formalizes is also tracked as Markdown at \texttt{docs/experiment-plan-predictor-stitching.md}.

\end{document}
